% Part: many-valued-logic
% Chapter: syntax-and-semantics
% Section: introduction

\documentclass[../../../include/open-logic-section]{subfiles}

\begin{document}

\olfileid{mvl}{syn}{int}

\olsection{పరిచయం}

సాంప్రదాయిక తర్కంలో $\lnot$, $\land$, $\lor$, $\lif$,
$\liff$ అనే ప్రతిజ్ఞావాక్య సంయోజకాలతో
\tetoken{ప్రతిజ్ఞావాక్య చరాల}{!!{propositional variable}s}
నుంచి నిర్మించిన \tetoken{సూత్రాలతో}{!!{formula}s}
వ్యవహరిస్తాం. దాని అర్థవిచారానికి $\True$, $\False$
అనే రెండు సత్యమూల్యాలను వాడతాం.
\tetoken{ప్రతిజ్ఞావాక్య చరాలను}{!!{propositional variable}s}
\tetoken{ఒక సత్యమూల్య కేటాయింపులో}{!!a{valuation}}~$\pAssign{v}$
అర్థనిర్దేశం చేస్తాం; అది \tetoken{ప్రతిజ్ఞావాక్య చరాలకు}{!!{propositional variable}s}
$\True$, $\False$ సత్యమూల్యాలను కేటాయిస్తుంది. అప్పుడు ఏ
\tetoken{సత్యమూల్య కేటాయింపు}{!!{valuation}} అయినా ప్రతి
\tetoken{సూత్రం}{!!{formula}}~$!A$కు సత్యమూల్యం
$\pValue{v}(!A)$ను నిర్ణయిస్తుంది. \tetoken{ఒక సూత్రం}{!!^a{formula}},
\tetoken{ఒక సత్యమూల్య కేటాయింపు}{!!a{valuation}}~$\pAssign{v}$లో
సంతృప్తమవడం, అంటే $\pSat{v}{!A}$ కావడం,
$\pValue{v}(!A) = \True$ అయినప్పుడూ, అప్పుడు మాత్రమే.

బహుమూల్య తర్కాలు సాంప్రదాయిక ద్విమూల్య తర్కానికి
సామాన్యీకరణలు: వాటిలో $\True$, $\False$లకే
పరిమితం కాకుండా మరిన్ని సత్యమూల్యాలను అనుమతిస్తాం.
కాబట్టి బహుమూల్య తర్కంలో \tetoken{ఒక సత్యమూల్య కేటాయింపు}{!!a{valuation}}~$\pAssign{v}$
అంటే, ప్రతి \tetoken{ప్రతిజ్ఞావాక్య చరం}{!!{propositional variable}}~$p$కు
సాధ్యమైన సత్యమూల్యాల పరిధి నుంచి ఒక విలువను కేటాయించే
ప్రమేయం. అనుమతించిన సత్యమూల్యాల సమితిని సాధారణంగా~$V$
అంటాం. సాంప్రదాయిక తర్కంలో $V = \{\True, \False\}$;
అక్కడ సంయోజకాల పరిచిత సత్య పట్టికలతో సూత్రపు
సత్యమూల్యం~$\pValue{v}(!A)$ను గణిస్తాం.

అదనపు సత్యమూల్యాలను చేర్చిన తరువాత, “మరియు,” “లేదా,”
“కాదు,” “అయితే---అప్పుడు” అని చదివే సంయోజకాలకు
$\pValue{v}(!A)$ను ఎలా గణించాలనే విషయంలో ఒకటి కంటే
ఎక్కువ సహజమైన ఎంపికలు ఉంటాయి. అందువల్ల బహుమూల్య
తర్కాన్ని సత్యమూల్యాల సమితి మాత్రమే నిర్ణయించదు;
ప్రతి సంయోజకానికి ఎంచుకున్న \emph{సత్యమూల్య ప్రమేయాలు}
కూడా నిర్ణయిస్తాయి. ఒక బహుమూల్య తర్కం~$\Log L$ కోసం
వాటిని ఎంచుకున్నాక, సాంప్రదాయిక తర్కంలోలాగే
సత్యమూల్యం $\pValue{v}(!A)[\Log L]$ను సత్యమూల్య
కేటాయింపే అనన్యంగా నిర్ణయిస్తుంది. సాంప్రదాయిక
తర్కంలాగే బహుమూల్య తర్కాలూ \emph{సత్యమూల్య-ప్రమేయాత్మకాలు}.

ఈ అర్థపర నిర్మాణ అంశాల ఆధారంగా సర్వసత్యం, అనుగమనం,
సంతృప్తిపరచదగినతనం అనే భావనలకు అనురూపాలను
నిర్వచించవచ్చు. సాంప్రదాయిక తర్కంలో
\tetoken{ఒక సూత్రం}{!!a{formula}} సర్వసత్యం కావడం,
ఏ~$\pAssign{v}$కైనా దాని సత్యమూల్యం
$\pValue{v}(!A) = \True$ అయినప్పుడూ, అప్పుడు మాత్రమే.
బహుమూల్య తర్కంలో దీన్నీ సామాన్యీకరించాలి. మొదటగా,
సత్యమూల్యాల సమితి~$V$లో $\True$ ఉండాలనే నియమం లేదు.
ఉదాహరణకు, కొన్ని బహుమూల్య తర్కాలు $0$ నుంచి~$1$
మధ్య ఉన్న అన్ని కరణీయ సంఖ్యలను సత్యమూల్యాలుగా
వాడతాయి. అటువంటి సందర్భంలో $1$ సాధారణంగా $\True$
పాత్ర పోషిస్తుంది. మరికొన్ని తర్కాల్లో ఒకటి కంటే
ఎక్కువ సత్యమూల్యాలు ఆ పాత్ర పోషిస్తాయి. కాబట్టి ప్రతి
బహుమూల్య తర్కానికి \emph{నిర్దేశిత సత్యమూల్యాల}
సమితి~$V^+$ ఉండాలని కోరుతాం. అప్పుడు
\tetoken{ఒక సూత్రం}{!!a{formula}},
\tetoken{ఒక సత్యమూల్య కేటాయింపు}{!!a{valuation}}~$\pAssign{v}$లో
సంతృప్తమవడం, అంటే $\pSat{v}{!A}[\Log L]$ కావడం,
$\pValue{v}(!A)[\Log L] \in V^+$ అయినప్పుడూ,
అప్పుడు మాత్రమే. \tetoken{ఒక సూత్రం}{!!^a{formula}}~$!A$
ఆ తర్కంలో సర్వసత్యం, $\Entails[\Log L] !A$ కావడం,
ఏ~$\pAssign{v}$కైనా $\pValue{v}(!A) \in V^+$
అయినప్పుడూ, అప్పుడు మాత్రమే. చివరగా,
\tetoken{సూత్రాల}{!!{formula}s} సమితి నుంచి~$!A$
అర్థపరంగా అనుగమిస్తుందని,
$\Gamma \Entails[\Log L] !A$ అని, $\Gamma$లోని అన్ని
\tetoken{సూత్రాలను}{!!{formula}s} సంతృప్తిపరిచే ప్రతి
\tetoken{సత్యమూల్య కేటాయింపు}{!!{valuation}} కూడా $!A$ను
సంతృప్తిపరిస్తే అంటాం.

\end{document}
